\workbookchapter{绪论}

\begin{exercises}
\begin{exercise} 在本章两步生成例子中，若将 $p(u\mid b)$ 改为 $0.7$，最大联合概率路径是否仍经过 $b$？写出全部路径概率后判断。

\begin{solution}
条件分布仍须归一化，故 $p(v\mid b)=0.3$。四条路径为 $0.6\times0.5=0.30$、$0.30$、$0.4\times0.7=0.28$、$0.4\times0.3=0.12$。最大路径改为经过 $a$ 的两条并列路径，不再经过 $b$。仅改变一个条件概率而不相应调整另一项，会破坏本题的概率模型。
\end{solution}
\end{exercise}

\begin{exercise} 同一段文本被两个分词器编码为不同长度，为什么不能仅比较每词元生成速度就断定哪个系统处理文本更快？

\begin{solution}
设同一文本编码长度为 $N_i$，每秒生成 $r_i$ 个词元，纯解码时间约为 $\frac{N_i}{r_i}$。较大的 $r_i$ 可能伴随更大的 $N_i$；还应计入分词、预填充和调度。应固定原始文本、生成质量与停止条件，比较完整请求时延、字符或字节处理量，并记录两套分词器和输入输出长度。
\end{solution}
\end{exercise}

\begin{exercise} 一个知识问答系统检索到了正确文档，却给出了错误答案。应检查哪些环节？为什么仅扩大模型不能保证解决问题？

\begin{solution}
依次检查正确片段是否进入实际上下文、是否被截断或错误序列化、来源版本与时间是否匹配、答案是否忠实使用证据，以及引用是否对应具体论断。针对故障可调整切分、重排、上下文编排、证据校验或拒答条件。扩大参数量不能修复丢失片段、错引用或权限过滤错误；修复后用同一问题及干扰变体回归。
\end{solution}
\end{exercise}

\begin{exercise} 某次修改使训练吞吐提高，但有效监督词元减少。需要补充哪些证据才能判断这是否是同一学习目标下的加速？

\begin{solution}
同时记录原始词元数、有效监督数、标签掩码、样本来源权重、批量及更新次数。若缩短序列或增加被屏蔽位置，原始词元每秒提高不代表有效学习加速。应保持目标位置和权重一致，比较有效监督词元每秒、单位更新损失及独立验证质量；若目标改变，明确作为新训练方案比较。
\end{solution}
\end{exercise}

\end{exercises}
